\documentclass[11pt,letterpaper]{article}

\usepackage[T1]{fontenc}
\usepackage{amsmath,amssymb}
\usepackage{tikz}
\usepackage{pgfplots}
\pgfplotsset{compat=1.18}
\usepackage{algorithm}
\usepackage{algpseudocode}
\usepackage{booktabs}
\usepackage{enumitem}
\usepackage[hidelinks]{hyperref}

\newcommand{\N}{\mathbb{N}}
\newcommand{\Z}{\mathbb{Z}}
\newcommand{\Q}{\mathbb{Q}}
\newcommand{\R}{\mathbb{R}}
\newcommand{\C}{\mathbb{C}}
\newcommand{\Map}[3]{#1 \colon #2 \to #3}
\newcommand{\Mod}[1]{#1\text{-}\mathrm{Mod}}
\newcommand{\GroupHom}{\mathrm{Hom}_{\mathbf{Grp}}}
\newcommand{\RingHom}{\mathrm{Hom}_{\mathbf{Ring}}}

\newtheorem{theorem}{Theorem}
\newtheorem{lemma}[theorem]{Lemma}
\newtheorem{corollary}[theorem]{Corollary}
\newtheorem{definition}[theorem]{Definition}
\newtheorem{remark}[theorem]{Remark}
\newenvironment{proof}{\noindent\textit{Proof.}}{\hfill\(\square\)}

\title{An ancestral inverse-odd join}
\author{ascendnoosphere}
\date{\today}

\begin{document}
\maketitle

\begin{abstract}
For integers \(a\ge 1\), \(b\ge 1\) and an odd integer \(u\), set
\(z=2^{a}3^{b}u-1\) and assume \(z\equiv 2^{a}-1\pmod{2^{a+2}}\).
The inverse odd step \(I(y)=(2y-1)/3\), repeated \(b\) times, is an exact
division by \(3\) at the first step and ends at \(2^{a+b}u-1\). One exact
division by \(2\) of the successor of that value produces
\(2^{a+b-1}u-1\). The congruence forces \(u\equiv(-1)^{b}\pmod{4}\).
Oddness of \(u\) blocks one further factor of \(2\). The congruence is a
hypothesis: the statement is not claimed for an arbitrary odd integer, and
it does not assert that the classical Collatz map reaches \(1\).
\end{abstract}

\section{The join}

Throughout, \(a\) and \(b\) are integers with \(a\ge 1\) and \(b\ge 1\), and
\(u\) is an odd integer. Set
\[
z=2^{a}3^{b}u-1,
\]
and assume the congruence
\begin{equation}
\label{eq:ancestral-congruence}
z\equiv 2^{a}-1\pmod{2^{a+2}}.
\end{equation}
All quotients below are exact integer divisions. The symbol \(I^{k}\) denotes
iteration of \(I\), with \(I^{0}\) the identity.

\begin{definition}
\label{def:ancestral-steps}
The inverse odd step and the successor half-shift are the integer operations
\[
I(y)=\frac{2y-1}{3},
\qquad
H(y)=\frac{y+1}{2}-1,
\]
defined when the displayed divisions are exact.
\end{definition}

\begin{lemma}
\label{lem:ancestral-closed}
Let \(z=2^{a}3^{b}u-1\). For every integer \(k\) with \(0\le k\le b\),
\[
I^{k}(z)=2^{a+k}3^{b-k}u-1,
\]
and each of the \(k\) divisions by \(3\) is exact.
\end{lemma}

\begin{proof}
The case \(k=0\) is the definition of \(z\). Suppose the formula holds for some
\(k\) with \(k+1\le b\), and write \(y=I^{k}(z)=2^{a+k}3^{b-k}u-1\).
Since \(k+1\le b\), the exponent \(b-k\) is at least \(1\), and
\[
b-k=\bigl(b-(k+1)\bigr)+1,
\qquad
a+(k+1)=(a+k)+1.
\]
Therefore
\begin{align*}
2y-1
&=2\bigl(2^{a+k}3^{b-k}u-1\bigr)-1 \\
&=2^{a+k+1}3^{b-k}u-2-1 \\
&=2^{a+k+1}3^{b-k}u-3 \\
&=3\bigl(2^{a+(k+1)}3^{b-(k+1)}u-1\bigr).
\end{align*}
The factor \(3\) divides \(2y-1\), and
\[
I(y)=2^{a+(k+1)}3^{b-(k+1)}u-1.
\]
Induction gives the formula at \(k=b\).
\end{proof}

\begin{lemma}
\label{lem:ancestral-first}
If \(b\ge 1\) and \(z=2^{a}3^{b}u-1\), then \(3\) divides \(2z-1\).
\end{lemma}

\begin{proof}
Since \(b\ge 1\), one has \(b=(b-1)+1\). Substitute
\(z=2^{a}3^{b}u-1\):
\[
2z-1=2\bigl(2^{a}3^{b-1}\cdot 3\cdot u-1\bigr)-1
=2^{a+1}3^{b-1}\cdot 3\cdot u-3
=3\bigl(2^{a+1}3^{b-1}u-1\bigr).
\]
The factor \(3\) divides \(2z-1\).
\end{proof}

\begin{lemma}
\label{lem:ancestral-half}
If \(n\ge 1\), then
\[
H\bigl(2^{n}u-1\bigr)=2^{n-1}u-1,
\]
and the division by \(2\) is exact.
\end{lemma}

\begin{proof}
The successor of \(2^{n}u-1\) is \(2^{n}u\). Since \(n\ge 1\),
\[
2^{n}u=2\cdot\bigl(2^{n-1}u\bigr),
\]
so the quotient by \(2\) is \(2^{n-1}u\). Subtracting \(1\) is the definition
of \(H\).
\end{proof}

\begin{lemma}
\label{lem:ancestral-odd-part}
If \(z=2^{a}3^{b}u-1\) and~\eqref{eq:ancestral-congruence} holds, then
\[
3^{b}u\equiv 1\pmod{4}.
\]
\end{lemma}

\begin{proof}
The difference of the two sides of~\eqref{eq:ancestral-congruence} is
\[
\bigl(2^{a}-1\bigr)-z=2^{a}\bigl(1-3^{b}u\bigr).
\]
The hypothesis says that \(2^{a+2}\) divides this difference. Since
\(2^{a+2}=2^{a}\cdot 4\) and \(2^{a}\neq 0\), it follows that \(4\) divides
\(1-3^{b}u\). That is the stated congruence.
\end{proof}

\begin{lemma}
\label{lem:ancestral-residue}
If \(3^{b}u\equiv 1\pmod{4}\), then \(u\equiv(-1)^{b}\pmod{4}\).
\end{lemma}

\begin{proof}
From \(3\equiv -1\pmod{4}\) one has \(3^{b}\equiv(-1)^{b}\pmod{4}\), and
therefore
\[
(-1)^{b}u\equiv 3^{b}u\equiv 1\pmod{4}.
\]
The square of the sign is \(1\):
\[
\bigl((-1)^{b}\bigr)^{2}=(-1)^{2b}=\bigl((-1)^{2}\bigr)^{b}=1.
\]
Multiplying the congruence by \((-1)^{b}\) yields
\[
u=\bigl((-1)^{b}\bigr)^{2}u\equiv(-1)^{b}\pmod{4}.
\]
\end{proof}

\begin{lemma}
\label{lem:ancestral-exact}
If \(a\ge 1\) and \(u\) is odd, then \(2^{a+b}\) does not divide
\(2^{a+b-1}u\).
\end{lemma}

\begin{proof}
Since \(a\ge 1\), one has \(a+b\ge 1\) and
\(2^{a+b}=2\cdot 2^{a+b-1}\). If \(2^{a+b}\) divided \(2^{a+b-1}u\),
cancellation of the nonzero factor \(2^{a+b-1}\) would give \(2\mid u\).
An odd integer is not divisible by \(2\).
\end{proof}

\begin{theorem}
\label{thm:ancestral-join}
Let \(a\ge 1\), \(b\ge 1\), and let \(u\) be an odd integer. Set
\(z=2^{a}3^{b}u-1\), and assume
\(z\equiv 2^{a}-1\pmod{2^{a+2}}\). Then \(3\) divides \(2z-1\),
\[
I^{b}(z)=2^{a+b}u-1,
\]
\[
H\bigl(I^{b}(z)\bigr)=2^{a+b-1}u-1,
\]
\(u\equiv(-1)^{b}\pmod{4}\), and \(2^{a+b}\) does not divide
\(2^{a+b-1}u\).
\end{theorem}

\begin{proof}
Lemma~\ref{lem:ancestral-first} is the divisibility by \(3\).
Lemma~\ref{lem:ancestral-closed} at \(k=b\) gives
\[
I^{b}(z)=2^{a+b}3^{0}u-1=2^{a+b}u-1.
\]
Since \(a\ge 1\), the exponent \(a+b\) is at least \(1\).
Lemma~\ref{lem:ancestral-half} therefore gives
\[
H\bigl(I^{b}(z)\bigr)=H\bigl(2^{a+b}u-1\bigr)=2^{a+b-1}u-1.
\]
Lemmas~\ref{lem:ancestral-odd-part} and~\ref{lem:ancestral-residue} give the
residue of \(u\). Lemma~\ref{lem:ancestral-exact} is the last clause.
\end{proof}

\begin{remark}
\label{rem:ancestral-scope}
The congruence~\eqref{eq:ancestral-congruence} is part of the hypothesis of
Theorem~\ref{thm:ancestral-join}. It is not asserted for every odd integer.
For \(a=1\), \(b=1\), and \(u=1\), one has \(z=5\) and
\(5\not\equiv 1\pmod{8}\), while \(u\not\equiv(-1)^{b}\pmod{4}\). That odd
integer lies outside the theorem. The argument does not show that the
classical Collatz map sends every positive integer to \(1\).
\end{remark}

\end{document}
